% =====================================================================
%  前置知识导引（译本 v2 新增，不属原书章节，不改变原书 1–10 章编号）
%  用途：把读懂本书所需的六组背景压缩成“定义 + 直觉 + 用在哪 + 文献”
%  体例：\chapter* + \section*，公式一律不编号，避免干扰原书公式编号
% =====================================================================

\chapter*{前置知识导引}
\markboth{前置知识导引}{前置知识导引}
\addcontentsline{toc}{chapter}{前置知识导引（译者增补）}

\noindent\textbf{说明。}本篇是中文译注本 v2 新增的导读，不属于原书，也不占用原书章节编号。原书自称“从零开始”，但它默认读者熟悉研究生代数量子化课程中的若干标准语言：李代数的根数据与泛包络代数、圈代数的中心扩张、仿射 Kac--Moody 代数、顶点代数、$D$-模与 Hecke 算子。缺少这些语言时，正文的许多“显然”会变成真正的障碍。本篇把六组背景各压缩成四件事：一句话定义、为什么这样定义（直觉）、在本书哪一章会真正用到、可进一步查阅的文献。页码均为原书 PDF 页码，方括号内文献沿用原书参考文献的缩写。

% =====================================================================
\section*{一、李代数与泛包络代数}
\addcontentsline{toc}{section}{一、李代数与泛包络代数}

\textbf{定义。}设 $\fg$ 是 $\C$ 上的有限维李代数，即带有反对称双线性括号 $[\ ,\ ]$ 且满足 Jacobi 恒等式的向量空间。若其 Killing 型 $(x,y)=\tr(\operatorname{ad} x\,\operatorname{ad} y)$ 非退化，则称 $\fg$ 半单；半单李代数同构于单李代数的直和。固定 Cartan 子代数 $\fh$，便有根空间分解
\[
  \fg=\fh\oplus\bigoplus_{\alpha\in\Delta}\fg_\alpha,\qquad
  [\fh,\fg_\alpha]\subset\fg_\alpha ,
\]
其中 $\Delta\subset\fh^*$ 是根集。取正根集 $\Delta_+$ 与单根 $\alpha_1,\dots,\alpha_r$，记 $\theta$ 为最高根，$\rho=\tfrac12\sum_{\alpha\in\Delta_+}\alpha$ 为 Weyl 向量，并令 $\alpha^\vee=2\alpha/(\alpha,\alpha)$。对偶 Coxeter 数定义为
\[
  h^\vee=1+\langle \rho,\theta^\vee\rangle ,
\]
于是 $\fsl_2$ 的 $h^\vee=2$，$\fsl_n$ 的 $h^\vee=n$。

泛包络代数 $U(\fg)$ 是张量代数模去关系 $xy-yx=[x,y]$ 所得的结合代数。PBW 定理说：若 $x_1,\dots,x_N$ 是 $\fg$ 的一组基，则单项式 $x_1^{a_1}\cdots x_N^{a_N}$（$a_i\ge 0$）构成 $U(\fg)$ 的基；等价地，取滤过之后 $\operatorname{gr}U(\fg)\cong\operatorname{Sym}(\fg)$。$U(\fg)$ 的中心 $Z(\fg)$ 在 $\fg$ 单半单时是 $r$ 元多项式环；二次 Casimir 元
\[
  \Omega=\sum_i x_i x^i
\]
（$x_i$ 与 $x^i$ 关于 Killing 型对偶）是其中最简单的元素，它在最高权为 $\lambda$ 的不可约模上作用为 $\tfrac12(\lambda,\lambda+2\rho)$。

\textbf{直觉。}$\fg$ 是群 $G$ 在单位元处的切空间，$U(\fg)$ 则是“支撑在单位元处的分布”或“非交换的多项式函数代数”。把 $G$ 的表示微分就得到 $\fg$ 的表示；反过来，单连通 $G$ 的每个有限维表示都来自 $\fg$ 的表示。本书大量使用 Verma 模：对 $\lambda\in\fh^*$，令 $\fb=\fh\oplus\bigoplus_{\alpha\in\Delta_+}\fg_\alpha$，则
\[
  M(\lambda)=U(\fg)\otimes_{U(\fb)}\C_{\lambda-\rho}
\]
是以 $\lambda$ 为最高权的 Verma 模；它有唯一不可约商 $L(\lambda)$。Harish-Chandra 定理把 $Z(\fg)$ 的特征标与 $\lambda$ 联系起来：$Z(\fg)$ 通过 $\chi_\lambda$ 作用在 $M(\lambda)$ 上，且 $\chi_\lambda=\chi_\mu$ 当且仅当 $\mu=w(\lambda+\rho)-\rho$ 对某个 外尔群元素 $w$ 成立。这套“最高权—Verma—中心特征标”的语言会在第 2、3、6、9 章反复出现。

\textbf{用在哪。}附录 A 给出系统复习；第 2 章用 $\fg$ 的根数据定义圈代数与仿射代数；第 6 章用 Verma 模与 Shapovalov 型；第 9 章用 $U(\fg)$ 的 Harish-Chandra 同态描述临界水平的中心。

\textbf{文献。}[Kac 1990] 第 1–2 章（根数据与 Kac--Moody 代数的一般框架）；[Dixmier 1977]（泛包络代数与 PBW）；[Humphreys 1972]（表示论入门，若手边有）。

% =====================================================================
\section*{二、中心扩张与 2-上循环}
\addcontentsline{toc}{section}{二、中心扩张与 2-上循环}

\textbf{定义。}李代数 $\mathfrak{a}$ 的中心扩张是短正合列
\[
  0\longrightarrow \C K\longrightarrow \hat{\mathfrak{a}}\longrightarrow \mathfrak{a}\longrightarrow 0
\]
且 $\C K$ 落在 $\hat{\mathfrak{a}}$ 的中心。若存在线性映射 $s:\mathfrak{a}\to\hat{\mathfrak{a}}$ 使 $s([x,y])-[s(x),s(y)]=\psi(x,y)K$，则 $\psi$ 是 $\mathfrak{a}$ 上的 2-上循环：
\[
  \psi([x,y],z)+\psi([y,z],x)+\psi([z,x],y)=0 .
\]
换一个截面 $s$ 只把 $\psi$ 改成 $\psi+\delta\beta$（上边界），所以中心扩张的同构类由 $H^2(\mathfrak{a};\C)$ 分类。对单李代数 $\fg$ 与 $R=\C((t))$，圈代数 $\fg((t))=\fg\otimes R$ 有
\[
  H^2(\fg((t));\C)\cong\C ,
\]
其生成元由不变双线性型 $(\ ,\ )$ 与留数给出：
\[
  \psi(x\otimes f,\,y\otimes g)=(x,y)\operatorname{Res}_{t=0}(f\,\dd g).
\]
于是
\[
  [x\otimes f,y\otimes g]=[x,y]\otimes fg+(x,y)\operatorname{Res}(f\,\dd g)\,K,\qquad [K,\hat{\mathfrak{g}}]=0 .
\]
若 $K$ 在表示上作用为标量 $k$，就称 $k$ 为该表示的\textbf{水平}。临界水平是
\[
  k=-h^\vee .
\]

\textbf{直觉。}中心扩张是“把射影表示变成真正的线性表示”的代数代价：圈群 $G((t))$ 的表示往往是射影的，扩张后的 $\hat{\mathfrak{g}}$ 把上循环吸收进交换子。物理上 $k$ 是 Wess--Zumino--Witten 模型的耦合常数，$k+h^\vee$ 是 Virasoro 中心荷的分子；当 $k=-h^\vee$ 时 Sugawara 构造的分母为零，Virasoro 作用退化，代数的中心却异常变大——这正是全书的出发点。

\textbf{用在哪。}第 1 章 §1.2 用中心扩张定义 $\hat{\mathfrak{g}}$；第 2 章 §2.2 计算上循环并引入临界水平；第 3 章说明临界水平中心的“异常大”；第 5、6 章用它解释 Wakimoto 模的构造障碍。

\textbf{文献。}[Kac 1990] 第 2 章与第 12 章；[Frenkel 2005b] 第 1 节；[Frenkel and Ben-Zvi 2004] 第 2 章。

% =====================================================================
\section*{三、仿射 Kac--Moody 代数}
\addcontentsline{toc}{section}{三、仿射 Kac--Moody 代数}

\textbf{定义。}在上一节的中心扩张上再加一个导子 $\dd$：
\[
  \hat{\mathfrak{g}}=\fg((t))\oplus\C K\oplus\C \dd,\qquad
  [\dd,x\otimes t^n]=n\,x\otimes t^n,\quad [\dd,K]=0 .
\]
这就是（未扭）仿射 Kac--Moody 代数。它的 Cartan 子代数是 $\hat{\fh}=\fh\oplus\C K\oplus\C \dd$，对偶空间中有虚根 $\delta$ 与基本权 $\Lambda_0$，满足
\[
  \langle \delta,K\rangle=0,\quad \langle\delta,\dd\rangle=1,\quad
  \langle\Lambda_0,K\rangle=1,\quad \langle\Lambda_0,\dd\rangle=0 .
\]
单根为 $\alpha_0=\delta-\theta,\alpha_1,\dots,\alpha_r$，单余根为 $\alpha_0^\vee=K-\theta^\vee,\alpha_1^\vee,\dots,\alpha_r^\vee$；外尔群是
\[
  W_{\mathrm{aff}}=W\ltimes Q^\vee ,
\]
其中 $Q^\vee$ 是余根格。对支配整权 $\Lambda$（即 $\langle\Lambda,\alpha_i^\vee\rangle\in\Z_{\ge0}$），不可约最高权模 $L(\Lambda)$ 是整水平的可积模；其形式特征标由 Weyl--Kac 公式给出。对一般 $\lambda\in\hat{\fh}^*$，Verma 模为
\[
  M(\lambda)=U(\hat{\mathfrak{g}})\otimes_{U(\hat{\fb})}\C_{\lambda-\rho},
\]
其中 $\hat{\fb}=\fb\oplus\C K\oplus\C \dd$。

\textbf{直觉。}仿射代数是“圈代数 + 中心 + 导子”的最小封闭系统：$\C K$ 记录射影上循环，$\dd$ 记录 $t$ 的次数（能量），两者一起把无限维的圈代数变成有根数据、有 外尔群、有最高权理论的 Kac--Moody 代数。水平 $k$ 决定中心荷；$k+h^\vee$ 出现在几乎所有公式的分母里，因此临界水平 $k=-h^\vee$ 是整本书的“奇点”，也是中心变大、opers 出现的地方。

\textbf{用在哪。}第 2 章用它定义真空模与 Verma 模；第 3 章研究其中心；第 6 章构造 Wakimoto 模；第 9 章讨论临界水平的 Verma/Weyl 模自同态。

\textbf{文献。}[Kac 1990] 第 6–7 章（根数据、Weyl--Kac 公式）；[Frenkel and Ben-Zvi 2004] 第 2 章；[Frenkel 2005b] 第 1–2 节。

% =====================================================================
\section*{四、顶点代数与 Wakimoto 模}
\addcontentsline{toc}{section}{四、顶点代数与 Wakimoto 模}

\textbf{定义。}顶点代数是一个四元组 $(V,|0\rangle,T,Y)$：$V$ 是向量空间，$|0\rangle\in V$ 是真空，$T:V\to V$ 是平移算子，$Y$ 是态—场对应
\[
  Y:V\longrightarrow \operatorname{End}(V)[[z,z^{-1}]],\qquad
  Y(a,z)=\sum_{n\in\Z}a_{(n)}z^{-n-1},
\]
满足真空公理 $Y(|0\rangle,z)=\id$、$Y(a,z)|0\rangle\in V[[z]]$ 且常数项为 $a$、平移公理 $[T,Y(a,z)]=\partial_zY(a,z)$，以及最重要的局域性公理：对任意 $a,b\in V$，存在 $N\gg0$ 使
\[
  (z-w)^N[Y(a,z),Y(b,w)]=0 .
\]
局域性等价于算子乘积展开（OPE）
\[
  Y(a,z)Y(b,w)\sim\sum_{n\ge0}\frac{Y(a_{(n)}b,w)}{(z-w)^{n+1}} .
\]
正规序乘积记为 $:Y(a,z)Y(b,w):$，即把 $a$ 的算子放在 $b$ 的左边。Virasoro 元 $\omega\in V$ 的场 $Y(\omega,z)=\sum L_nz^{-n-2}$ 给出中心荷为 $c$ 的 Virasoro 代数。

对仿射代数 $\hat{\mathfrak{g}}$，真空模
\[
  V^k(\fg)=U(\hat{\mathfrak{g}})\otimes_{U(\fg[[t]]\oplus\C K)}\C_k
\]
（$\fg[[t]]$ 通过求值作用，$K$ 作用为 $k$）是一个顶点代数，称为通用仿射顶点代数；其唯一单商记为 $L_k(\fg)$。当 $k$ 不是临界水平时，Sugawara 构造给出 Virasoro 元；当 $k=-h^\vee$ 时分母为零，Virasoro 元不再存在，而中心 $Z(\hat{\mathfrak{g}})$ 反而变大。

Wakimoto 模是 $V^k(\fg)$ 的自由场实现：对 $\fsl_2$，取一对权重 $(1,0)$ 的自由场 $\beta,\gamma$（$\beta$ 为玻色鬼场、$\gamma$ 为其共轭）以及一个自由玻色子 $\varphi$，用它们构造 $\hat{\fsl}_2$ 的作用，使真空模等同于 Fock 空间；对一般 $\fg$，取主幂零元 $e\in\fg$ 对应的 $\beta\gamma$ 系统与 Cartan 部分的自由玻色子。这样得到的模在临界水平给出中心作用的显式计算，也是证明 Feigin--Frenkel 同构的关键工具。屏蔽算子（screening operator）是形如
\[
  Q_i=\oint \ee^{\alpha_i\cdot\varphi(z)}\,\dd z
\]
的顶点算子，它与 $\hat{\mathfrak{g}}$ 的作用交换，用来切出可积模或单商。

\textbf{直觉。}顶点代数把“两个算子在同一点相乘”的奇性编码成代数结构：OPE 的极点阶数就是算子乘积展开的阶数。自由场实现则是“用交换的玻色子模拟非交换的仿射代数”，代价是要引入鬼场并证明上同调消去；Wakimoto 模正是这样一个实现，它让原本难以计算的中心作用变成自由场上的微分算子。

\textbf{用在哪。}第 2 章 §2.4 给出顶点代数速成；第 3 章用顶点代数语言重述中心；第 5–8 章用 Wakimoto 模与屏蔽算子完成主定理；第 9 章用其计算 Verma 模自同态。

\textbf{文献。}[Frenkel and Ben-Zvi 2004]（顶点代数与代数曲线的标准参考）；[Kac 1998]（入门）；[Frenkel 2005b] 第 3–6 节（Wakimoto 模与屏蔽算子）；[Wakimoto 1986]（原始构造）。

% =====================================================================
\section*{五、$D$-模与本征层}
\addcontentsline{toc}{section}{五、$D$-模与本征层}

\textbf{定义。}设 $X$ 是光滑复代数曲线。$D_X$ 是 $X$ 上微分算子构成的层；左 $D_X$-模的范畴记为 $D(X)\text{-}\mathrm{mod}$。当模还是 $\cO_X$-凝聚时，左 $D_X$-模等价于带平坦联络的向量丛：联络 $\nabla$ 满足 $\nabla_{f s}=f\nabla s+(\dd f)s$ 且曲率为零。Riemann--Hilbert 对应说，在 $X$ 为紧黎曼面时，$D_X$-模的导出范畴与 $\C$ 上有限维局部系统的导出范畴等价（正则奇点情形要加上 Stokes 结构）。

本书更关心的是模空间上的 $D$-模。设 $G$ 是复约化群，$\Bun_G(X)$ 是 $X$ 上 $G$-主丛的模空间；$D(\Bun_G(X))\text{-}\mathrm{mod}$ 是几何 Langlands 对应的一侧。另一侧是 $\Loc_{{}^L G}(X)$，即 $X$ 上 ${}^L G$-局部系统的模空间。几何 Langlands 猜想（无分歧情形）断言，存在 $D(\Bun_G(X))\text{-}\mathrm{mod}$ 与 $\cO(\Loc_{{}^L G}(X))\text{-}\mathrm{mod}$ 之间的等价，且把每个局部系统 $E$ 送到以 $E$ 为特征的 Hecke 本征层。2007 年时这仍主要是猜想；本书第 10 章讨论的是它的局部版本。

\textbf{直觉。}$D$-模是“带联络的向量丛”的范畴化：$\cO_X$-模只看函数，$D_X$-模还看沿向量场的求导。把它搬到 $\Bun_G$ 上，Hecke 算子就成为“在一点处修改丛”的对应，而几何 Langlands 说这些修改算子的本征对象恰好由 ${}^L G$-局部系统参数化。

\textbf{用在哪。}第 1 章 §1.2–§1.4 用局部系统与 $D$-模重述局部参数；第 4 章用 $D^\times$-上的 $D$-模定义 opers；第 9、10 章用 $\Bun_G$、$\Gr$ 上的 $D$-模表述局部几何 Langlands。

\textbf{文献。}[Beilinson and Drinfeld 1997]（几何 Langlands 的纲领性文献）；[Frenkel 2007]（综述）；[Borel et al. 1987]（$D$-模入门）。

% =====================================================================
\section*{六、Hecke 算子与几何 Langlands}
\addcontentsline{toc}{section}{六、Hecke 算子与几何 Langlands}

\textbf{定义。}固定 $X$ 上一点 $x$。Hecke 对应是 $\Bun_G(X)$ 上的一个对应：
\[
  \Bun_G(X)\xleftarrow{\ p\ }\mathcal{H}_x\xrightarrow{\ q\ }\Bun_G(X),
\]
其中 $\mathcal{H}_x$ 参数化“在 $x$ 之外同构、在 $x$ 处允许一个 $G[[t]]$-修改”的丛对 $(E,E')$。对 ${}^L G$ 的每个有限维表示 $V$，Hecke 算子
\[
  \mathrm{H}_{V,x}:D(\Bun_G(X))\text{-}\mathrm{mod}\longrightarrow D(\Bun_G(X))\text{-}\mathrm{mod}
\]
由 $\mathrm{H}_{V,x}(\cF)=q_*(\,p^*\cF\otimes\cV\,)$ 给出，其中 $\cV$ 是 $\Gr_G$ 上由 $V$ 经几何 Satake 对应得到的 perverse 层。一个 $D$-模 $\cF$ 称为 \textbf{Hecke 本征层}，若对每个 $V$ 都有
\[
  \mathrm{H}_{V,x}(\cF)\cong \cF\boxtimes V
\]
（在 $x$ 的附近的意义下），即 $\cF$ 是所有 Hecke 算子的共同本征对象。

局部情形把 $\Bun_G(X)$ 换成仿射 Grassmannian $\Gr_G=G((t))/G[[t]]$，把 Hecke 算子换成卷积；球面 Hecke 代数 $\mathcal{H}(G((t)),G[[t]])$ 通过 Satake 同构同构于 $\Rep({}^L G)$，其几何版本是 Mirković--Vilonen 的几何 Satake 等价：
\[
  \Rep({}^L G)\;\simeq\;\operatorname{Perv}_{G[[t]]}(\Gr_G).
\]

\textbf{直觉。}Hecke 算子是“在一点插入一个表示”的算子；Hecke 本征层是它的本征向量。几何 Langlands 的整个内容可以概括为：对每个 ${}^L G$-局部系统 $E$，存在唯一的 Hecke 本征层 $\cF_E$，其本征值由 $E$ 在 $x$ 处的单值性给出。Satake 同构说明“插入哪个表示”由 ${}^L G$ 的表示论分类，这正是 Langlands 对偶群出现在几何侧的原因。

\textbf{用在哪。}第 1 章 §1.3–§1.5 用 Hecke 本征层陈述无分歧几何 Langlands；第 9 章用局部 Hecke 作用定义 $\hat{\mathfrak{g}}$-模的局部化；第 10 章把 Hecke 算子与 opers 联系起来。

\textbf{文献。}[Beilinson and Drinfeld 1997]；[Frenkel 2007]；[Frenkel et al. 2002]（几何 Langlands 的构造）；[Mirković and Vilonen 2007]（几何 Satake，若手边有）。

\bigskip
\noindent\textbf{阅读建议。}第一次读本书时，六节中至少要先掌握第二、三节（中心扩张与仿射代数）和第四节前半（顶点代数的定义与真空模），否则第 2 章会读得很慢。第五节与第六节可以等到第 9、10 章再回头补；这两节的作用是让“几何 Langlands”这句话在读者心里有一个可操作的含义，而不是一句口号。
